\documentclass[11pt,a4paper]{article}

\usepackage[T1]{fontenc}
\usepackage[utf8]{inputenc}
\usepackage[brazil]{babel}
\usepackage{lmodern}
\usepackage{microtype}
\usepackage[a4paper,margin=2.5cm,headheight=15pt]{geometry}
\usepackage{array,booktabs,longtable,tabularx}
\usepackage{enumitem}
\usepackage{xcolor}
\usepackage{hyperref}
\usepackage{fancyhdr}
\usepackage{titlesec}

\definecolor{upeblue}{HTML}{123B62}
\definecolor{lightblue}{HTML}{EAF1F7}
\definecolor{darkgray}{HTML}{3A3A3A}

\hypersetup{
  colorlinks=true,
  linkcolor=upeblue,
  urlcolor=upeblue,
  pdftitle={Especificação B1 - PlantUML Parser HSL},
  pdfauthor={Helaine Solange Lins Barreiros},
  pdfsubject={Registro de Programa de Computador}
}

\pagestyle{fancy}
\fancyhf{}
\lhead{\small Especificação B1 - Parser HSL}
\rhead{\small Registro de Programa de Computador}
\cfoot{\small Página \thepage}

\titleformat{\section}{\Large\bfseries\color{upeblue}}{\thesection}{0.7em}{}
\titleformat{\subsection}{\large\bfseries\color{darkgray}}{\thesubsection}{0.7em}{}
\setlist[itemize]{leftmargin=1.5em,itemsep=2pt,topsep=4pt}
\setlist[enumerate]{leftmargin=1.7em,itemsep=2pt,topsep=4pt}
\setlength{\parindent}{0pt}
\setlength{\parskip}{6pt}
\setlength{\emergencystretch}{2em}
\newcommand{\code}[1]{\texttt{\detokenize{#1}}}
\newcommand{\pending}{\textbf{PENDENTE}}

\begin{document}
\pagenumbering{gobble}

\begin{titlepage}
  \centering
  \vspace*{1.3cm}
  {\large\bfseries ESPECIFICAÇÃO B1\par}
  \vspace{0.25cm}
  {\large Registro de Programa de Computador\par}
  \vspace{2.3cm}
  {\Huge\bfseries\color{upeblue} PlantUML Parser HSL com Saída Estruturada para Casos de Uso\par}
  \vspace{0.8cm}
  {\Large Produto 4 - Versão candidata 1.0.0\par}
  \vfill
  \begin{tabular}{>{\bfseries}p{4.0cm}p{9.0cm}}
    Autora das modificações & Helaine Solange Lins Barreiros \\
    Obra originária & \code{Enteee/plantuml-parser} \\
    Licença & Apache License 2.0 \\
    Commit-base & \code{966a72e} \\
    Commit funcional & \code{f317b79} \\
    Branch & \code{codex/registro-parser-usecase-v1.0.0} \\
    Estado & Minuta técnica completa; pacote e SHA-256 ainda não gerados \\
  \end{tabular}
  \vfill
  {\large 30 de outubro de 2026\par}
  \vspace{0.8cm}
  \fcolorbox{upeblue}{lightblue}{\parbox{0.88\textwidth}{\small
  O objeto é a versão derivada que acrescenta suporte estruturado a diagramas de casos de uso. O pedido não reivindica autoria ou exclusividade sobre o código preexistente. A licença, a obra originária e o escopo das modificações devem permanecer identificados.}}
\end{titlepage}

\pagenumbering{roman}
\tableofcontents
\clearpage
\pagenumbering{arabic}

\section{Introdução}

\subsection{Finalidade}

O programa analisa texto PlantUML e fornece uma representação estruturada, em JSON, de diagramas de casos de uso. A saída separa atores, casos de uso, pacotes, retângulos e relações, preservando atributos relevantes para processamento automatizado por outros programas.

\subsection{Contexto de desenvolvimento}

A extensão foi desenvolvida para apoiar a avaliação programática de diagramas UML no doutorado. O parser originário já realizava parsing de PlantUML; a versão HSL modifica e amplia essa base para reconhecer construções necessárias aos casos de uso e emitir uma saída organizada para consumo pelo motor de avaliação.

\subsection{Natureza derivada}

A versão conserva arquitetura, gramática e código do projeto originário Enteee/\allowbreak plantuml-parser, licenciado sob Apache-2.0, e modifica essa base. Por isso, tecnicamente é obra derivada. Essa qualificação descreve a proveniência e não impede que as modificações autorais formem um produto relacionado e um pedido próprio, desde que a origem e a licença sejam respeitadas.

\section{Identificação}

\begin{longtable}{>{\bfseries}p{4.1cm}p{10.0cm}}
\toprule
Campo & Informação \\
\midrule
\endhead
Nome do programa & PlantUML Parser HSL com Saída Estruturada para Casos de Uso \\
Nome curto & PlantUML Parser HSL - Use Case Structured \\
Versão & v1.0.0 candidata; confirmar após o congelamento \\
Data de criação & Pendente de validação; marco funcional no commit \code{f317b79} \\
Data de publicação & Pendente de orientação do NIT; não inferir automaticamente do GitHub \\
Linguagens & TypeScript, JavaScript e PEG.js \\
Ambiente & Biblioteca Node.js e executável de linha de comando \\
Natureza & Versão derivada de software livre, com modificações autorais identificadas \\
Licença & Apache-2.0, preservada da obra originária \\
\bottomrule
\end{longtable}

\section{Problema, finalidade e usuários}

\subsection{Problema atendido}

A árvore sintática genérica de um parser não fornece necessariamente um contrato direto para avaliação de diagramas de casos de uso. O software acrescenta reconhecimento e formatação capazes de produzir categorias explícitas e relacionamentos estruturados, reduzindo o acoplamento dos consumidores à estrutura interna da árvore de parsing.

\subsection{Finalidade}

\begin{itemize}
  \item reconhecer elementos de diagramas de casos de uso em PlantUML;
  \item representar atores e seus estereótipos;
  \item representar casos de uso, pacotes e retângulos;
  \item preservar a associação dos elementos aos grupos;
  \item estruturar relações, extremidades, corpos de seta, rótulos e tipo de relação;
  \item oferecer formatadores JSON consumíveis por ferramentas externas;
  \item disponibilizar a variante pelo módulo e pelo executável \code{plantuml-parser-hsl}.
\end{itemize}

\subsection{Usuários-alvo}

Desenvolvedores, pesquisadores e ferramentas que precisam transformar diagramas de casos de uso PlantUML em dados estruturados para validação, comparação, métricas ou integração.

\section{Fronteira do produto}

\subsection{Escopo incluído nesta versão}

\begin{itemize}
  \item alterações de gramática para construções utilizadas em casos de uso;
  \item ampliação dos tipos internos para atores, casos de uso, grupos e relações;
  \item formatadores \code{structured} e \code{usecase_structured};
  \item integração dos formatadores à API exportada;
  \item ajustes do formatador de grafo relacionados aos novos tipos;
  \item configuração de compilação, manifesto NPM e lockfile compatíveis;
  \item exemplo de diagrama e script de verificação existentes no commit da extensão.
\end{itemize}

\subsection{Escopo explicitamente excluído}

\begin{itemize}
  \item extensões posteriores para diagramas de classe, sequência ou atividade;
  \item alterações posteriores aos commits \code{f317b79} e \code{8446cf4};
  \item Motor de Avaliação Técnica e Estatística em Python;
  \item resultados, diagramas e planilhas empíricas do doutorado;
  \item \code{node_modules}, arquivos compilados, cobertura, caches e tokens;
  \item qualquer reivindicação sobre o código originário não modificado.
\end{itemize}

\subsection{Arquivos modificados ou acrescentados desde o commit-base}

O diff \code{966a72e..f317b79} identifica o núcleo da versão:

\begin{itemize}
  \item \code{src/formatters/graph.ts}: ajustes de atores e relações;
  \item \code{src/formatters/structured.ts}: novo formatador estruturado;
  \item \code{src/formatters/usecase_structured.ts}: formatador especializado;
  \item \code{src/index.ts}: exportação e registro dos formatadores;
  \item \code{src/plantuml.pegjs}: extensões da gramática;
  \item \code{src/types.ts}: tipos para casos de uso, atores, grupos e relações;
  \item \code{package.json}, \code{package-lock.json} e \code{tsconfig.json}: manifesto, dependências e compilação;
  \item \code{actor_test.plantuml} e \code{test-actor.js}: entrada e verificação de apoio.
\end{itemize}

\section{Funcionalidades principais}

\begin{longtable}{p{0.8cm}p{7.4cm}p{5.6cm}}
\toprule
ID & Funcionalidade & Evidência principal \\
\midrule
\endhead
F1 & Analisar fonte PlantUML usando a gramática estendida & \code{src/plantuml.pegjs} \\
F2 & Representar atores, casos de uso e grupos por tipos próprios & \code{src/types.ts} \\
F3 & Extrair atores e estereótipos & \code{usecase_structured.ts} \\
F4 & Extrair casos de uso & \code{usecase_structured.ts} \\
F5 & Organizar elementos contidos em pacotes e retângulos & \code{usecase_structured.ts} \\
F6 & Estruturar relações e detalhes das setas & \code{usecase_structured.ts} \\
F7 & Inferir tipos conhecidos das extremidades das relações & \code{usecase_structured.ts} \\
F8 & Emitir JSON organizado por categorias & formatadores estruturados \\
F9 & Expor os novos formatadores na biblioteca & \code{src/index.ts} \\
F10 & Disponibilizar execução pelo binário NPM & \code{package.json} \\
\bottomrule
\end{longtable}

\section{Arquitetura e contratos}

\begin{center}
\fcolorbox{upeblue}{lightblue}{\parbox{0.91\textwidth}{\centering
Arquivo ou texto PlantUML\\[4pt]
$\Downarrow$\\[4pt]
Gramática PEG.js $\longrightarrow$ objetos tipados e árvore de parsing\\[4pt]
$\Downarrow$\\[4pt]
Formatador \code{usecase_structured}\\[4pt]
$\Downarrow$\\[4pt]
JSON: atores, casos de uso, pacotes, retângulos e relações\\[4pt]
$\Downarrow$\\[4pt]
Programas consumidores}}
\end{center}

\subsection{Entradas}

\begin{itemize}
  \item texto PlantUML ou caminho/glob de arquivos PlantUML;
  \item opções já suportadas pela API de parsing;
  \item seleção do formatador de saída.
\end{itemize}

\subsection{Saídas}

O JSON contém:

\begin{itemize}
  \item \code{actors}: nome, título e estereótipos;
  \item \code{usecases}: nome e título;
  \item \code{packages} e \code{rectangles}: identificação e elementos internos;
  \item \code{relationships}: origem, destino, tipos conhecidos, tipo de relação, rótulo e partes das setas.
\end{itemize}

\subsection{Relação com o Produto 3}

O Produto 3 chama o executável \code{plantuml-parser-hsl} por meio de um adaptador Python e consome o JSON. Essa relação é de interoperabilidade: o código TypeScript do parser permanece neste repositório e não precisa ser copiado integralmente para o pacote do Produto 3.

\section{Proveniência e autoria}

\subsection{Obra originária}

\begin{itemize}
  \item projeto: \code{Enteee/plantuml-parser};
  \item repositório: \url{https://github.com/Enteee/plantuml-parser};
  \item autor originário indicado no manifesto: Ente;
  \item licença: Apache License 2.0;
  \item commit-base adotado: \code{966a72e}.
\end{itemize}

\subsection{Modificações desta versão}

O histórico entre \code{966a72e} e \code{f317b79} contém um commit funcional atribuído a Helaine Barreiros. A documentação de saneamento foi acrescentada no commit \code{8446cf4}. O pedido deve identificar Helaine como autora das modificações desta versão, sem substituir os créditos ou direitos sobre a obra originária.

\subsection{Contribuições posteriores fora do objeto}

As extensões de classes, sequências e atividades existentes em outras branches ou commits posteriores não pertencem a esta versão. Seus autores não devem ser incluídos automaticamente neste pedido, salvo decisão de ampliar a fronteira e gerar novo congelamento.

\section{Originalidade e alcance dos direitos}

A originalidade candidata está nas modificações da gramática e dos tipos e na implementação dos formatadores que organizam elementos de casos de uso em contrato JSON específico. O alcance do pedido deve se limitar às modificações e à seleção ou organização autoral, respeitando os direitos e avisos da base Apache-2.0.

A licença autoriza uso, modificação e redistribuição sob suas condições, mas sua suficiência documental para o processo brasileiro deve ser confirmada pelo NIT. O arquivo \code{LICENSE} deve acompanhar o pacote, e avisos de copyright existentes não devem ser removidos.

\section{Classificação candidata}

\begin{longtable}{>{\bfseries}p{4.0cm}p{10.0cm}}
\toprule
Categoria & Opções candidatas \\
\midrule
\endhead
Campo de aplicação & CO03 - linguagem artificial; IF02 - processamento de informação; IF10 - processamento de dados \\
Tipo de programa & DS01 - ferramenta de suporte ao desenvolvimento; DS05 - biblioteca de rotinas; DS06 - apoio à programação; TC01 - aplicações técnico-científicas \\
Linguagens & TypeScript, JavaScript e PEG.js \\
Forma de apresentação & código-fonte e documentação técnica identificados por SHA-256 \\
\bottomrule
\end{longtable}

As opções são propostas técnicas e devem ser confirmadas nas tabelas vigentes do INPI e pelo NIT.

\section{Privacidade e configuração externa}

A auditoria da branch não identificou CPF, telefone, endereço pessoal, caminho local, chave privada ou valor de serviço privado no código desta versão. Referências a \code{secrets} em workflows representam nomes de segredos, não seus valores.

\begin{itemize}
  \item tokens NPM e GitHub devem permanecer nos mecanismos de segredo das plataformas;
  \item \code{.env}, configurações locais e credenciais não integram o pacote;
  \item URLs da obra originária e metadados de licença permanecem como proveniência;
  \item o ZIP final deve ser novamente verificado antes do cálculo do hash.
\end{itemize}

\section{Pacote depositável e forma de entrega}

\subsection{Conteúdo a incluir}

\begin{itemize}
  \item arquivos do núcleo listados na Seção 4.3;
  \item arquivos necessários para compilar a gramática e o TypeScript;
  \item \code{package.json} e lockfile;
  \item testes e exemplos sintéticos pertinentes;
  \item \code{LICENSE} e demais avisos aplicáveis;
  \item README com instalação, entrada, saída e exemplo;
  \item esta especificação e manifesto de arquivos;
  \item relatório ou patch que documente \code{966a72e..f317b79}.
\end{itemize}

\subsection{Conteúdo a excluir}

\begin{itemize}
  \item \code{.git} e branches não pertencentes à versão;
  \item \code{node_modules}, \code{dist}, \code{build}, cobertura e caches;
  \item tokens, \code{.env} e configurações locais;
  \item arquivos de classes, sequências ou atividades posteriores;
  \item dados e resultados empíricos;
  \item código do Produto 3 ou de outros produtos.
\end{itemize}

\subsection{Artefatos de entrega}

\begin{enumerate}
  \item commit ou tag da versão congelada;
  \item ZIP determinístico;
  \item manifesto de arquivos e relatório de proveniência;
  \item SHA-256 do ZIP;
  \item licença e avisos;
  \item esta especificação e instruções de reconstrução;
  \item cópia exata do ZIP em ao menos dois locais controlados.
\end{enumerate}

\section{Verificação e limitações}

\begin{longtable}{p{3.4cm}p{5.3cm}p{5.3cm}}
\toprule
Item & Estado & Ação antes do hash \\
\midrule
\endhead
Fronteira autoral & delimitada por \code{966a72e..f317b79} & gerar relatório final do diff \\
Dados pessoais e segredos & não encontrados no snapshot auditado & repetir varredura no ZIP \\
Dependências Node & declaradas nos manifestos & reinstalar de forma limpa \\
Compilação e testes & não comprovados no ambiente final & executar \code{npm ci}, build e testes \\
Exemplo de caso de uso & existe entrada de apoio & transformar em teste reproduzível \\
Licença e proveniência & Apache-2.0 e upstream identificados & validar documentos com o NIT \\
\bottomrule
\end{longtable}

\section{Resumo digital}

\begin{longtable}{>{\bfseries}p{4.0cm}p{10.0cm}}
\toprule
Campo & Valor \\
\midrule
\endhead
Algoritmo & SHA-256 \\
Nome sugerido & \code{produto4-plantuml-parser-hsl-usecase-v1.0.0.zip} \\
Resumo digital & \pending{} - gerar após manifesto, testes e congelamento \\
Commit ou tag do pacote & \pending \\
Custódia do arquivo & \pending{} DE DEFINIÇÃO INSTITUCIONAL \\
\bottomrule
\end{longtable}

\section{Pendências institucionais}

\begin{itemize}
  \item confirmar autoria formal das modificações e titularidade correspondente;
  \item confirmar como a obra originária ou seus titulares devem constar no pedido;
  \item validar a Apache-2.0 como documento de autorização e os anexos necessários;
  \item decidir data de criação e publicação;
  \item confirmar classificações do INPI;
  \item aprovar a fronteira restrita a casos de uso;
  \item gerar o pacote final e seu SHA-256.
\end{itemize}

\end{document}
